% RSE: Section 6, Conclusions (under 250 words; the "how much / where" contrast used once).
Dated Landsat waterlines, levelled by satellite altimetry, record where the North Caspian shoreline has been since the start of the altimetry record. Tested against them, four of five global bathymetry grids misplace the shoreline in nearly every sector, even after a vertical offset fitted to flooded area. On the Kizlyar--Dagestan and Volga--Ural interfluve coasts, where independent ICESat-2 heights confirm the observations, the misses are \metricdp{rev.s.headline.west.d_km_min}{1}--\metricdp{rev.s.headline.west.d_km_max}{1}~km, several times the observation noise floor. Whole-metre storage, coarse cells, level errors and datum choices do not explain them. The offset the grids need changes with level, so no single shift fixes them. ETOPO~2022 passes only where its shoreline comes from land heights measured after the seabed emerged. Water-level pathways decide how much seabed emerges by mid-century; the grids cannot say where. Emergence and exposure assessments should start from the latest observed shoreline and test any grid against dated shorelines before use.
